\documentclass{article}

\usepackage{amsmath}
\usepackage{bm}

\begin{document}
This is to build up some intuition about option pricing and trading
strategies. We will use the following notation:
\begin{itemize}
    \item $S_0$: Current stock price
    \item $K$: Strike price of option
    \item $T$: Time to expiration of option
    \item $S_T$: Stock price at maturity
    \item $r$: Continuously compounded risk-free rate of interest for an investment maturing in time $T$
    \item $C$: Value of American call option to buy on share
    \item $P$: Value of American put option to sell one share
    \item $c$: Value of European call option to buy one share
    \item $p$: Value of European put option to sell one share
    \item $D$: Present value of dividends paid during the life time of the option
\end{itemize}

\section{Interesting Portfolios}
Consider the following 4 portfolios for a non-dividend-paying stock
\begin{itemize}
    \item Portfolio A: one European call option plus a mount of cash equal to \(Ke^{-rT}\)
    \item Portfolio B: one European put option plus one share	    
    \item Portfolio C: one share
    \item Portfolio D: an amount of cash equal to \(Ke^{-rt}\)
\end{itemize}
These 4 portfolios should be very reasonable and intuitive to think about. 
Portfolio A holds a call with potential cash to buy stock. 
Portfolio B holds a put with a share to potentially sell. 
Portfolio C holds a single share while portfolio D holds just cash.
At time $T$, these 4 portfolios are:
\begin{itemize}
    \item Portfolio A: 
    if \(S_T \le K\), option worths nothing, the cash equals to $K$.
    if \(S_T > K\), we exercise the option by using the cash that 
    is equal to $K$ to to buy one share which is valued at $S_T$. 
    So the value of the portfolio is \(\max(S_T, K)\)
    \item Portfolio B: 
    if \(S_T \ge K\), option worths nothing, we have one share that is
    valued at $S_T$.
    if \(S_T < K\), we exercise the option by selling the share at $K$ and 
    we hold cash equal to $K$ at the end. 
    So the value of the portfolio is \(max(S_T, K)\). 
    \item Portfolio C: one share
    \item Portfolio D: an amount of cashe equal to $K$
\end{itemize}

This can be summarized as the following table content:

\begin{tabular}{c c c c}
    \hline \\
    Portfolio & t=0 & t=T & Note \\
    \hline \\
    A & \(c + Ke^{-rT}\) & \(\max(S_T, K)\) & long call with cash \\
    B & \(p + S_0\) & \(\max(S_T, K)\) & protected put \\
    C & $S_0$ & $S_T$ & underlying stock \\
    D & \(Ke^{-rT}\) & $K$ & cash \\
    \hline
\end{tabular}

In scenarios the stock pays future dividend, denote $D$ as the present 
value of the dividends during the life of the option. We would have the
similar 4 portfolios:
\begin{itemize}
    \item Portfolio A': one European call option plus a mount of cash equal to \(D+Ke^{-rT}\)
    \item Portfolio B': one European put option plus one share	    
    \item Portfolio C': one share
    \item Portfolio D': an amount of cash equal to \(D+Ke^{-rt}\)
\end{itemize}

\section{Basic Boundaries}
\begin{itemize}
    \item American options are more valuable than European options: 
    \(C \ge c, P \ge p\)
    \item call options can't be worth more than the stock: 
    \(C \le S_0, c \le S_0\)
    \item put options gives the right to sell one share for $K$, 
    so it can't be worth more than $K$: 
    \(p \le K, P \le K\)	    
\end{itemize}

\subsection{Boundaries for European call options}
Consider Portfolio A and C. At time $T$, you can observe that the
value of Portfolio A always worth as much as, and can be worth more than
Portfolio C. \textit{It follows that in the absence of arbitrage opportunites this
must also be true today}. Hence,
\[
c + Ke^{-rT} \ge S_0
\]
Combine with the fact that option price needs to be positve and the boundary 
established above, we have:
\[
\max(S_0 - Ke^{-rT}, 0) \le c \le S_0
\]
Similarly, for a stock that pays dividends, we have:
\[
\max(S_0 - D - Ke^{-rT}, 0) \le c
\]

\subsection{Boundaries for European put options}
Consider Portfolio B and D. At time $T$, you can observe that the value 
of Portfolio B always worth as much as, and can be worth more than
Portfolio D. \textit{It follows that in the absense of arbitrage opportunites
this must also be true today}. Hence,
\[
p + S_0 \ge Ke^{-rT}
\]
Combine with the fact that option price needs to be positve and the boundary
established earlier, we have:
\[
\max(Ke^{-rT} - S_0, 0) \le p \le Ke^{-rT}
\]
Similarly, for a stock that pays dividends, we have:
\[
\max(D + Ke^{-rT} - S_0, 0) \le p
\]

\section{Pull-Call Partity}
\subsection{European options}
Consider Portfolio A and B. At time $T$ they always worth the same. Because
they can't be exercised before the expiration at time $T$, the portfolios must
therefore have identical values today. Hence
\[
c + Ke^{-rT} = p + S_0
\]
Similarly, for a stock that pays dividends, we have:
\[
c + D + Ke^{-rT} = p + S_0
\]

\subsection{American options}
\subsection{Trading Strategies with a single option}
Let's consider trading strategies that can be explained by put-call parity.
\[
p + S_0 = c + Ke^{-rT} + D
\]
This can be read as: a long position in a put with a long position in the stock is
equivalent to a long call position plus a certain amount of cash $Ke^{-rT} + D$.
This is a \textbf{protected put}.

Similarly consider this form of put-call parity:
\[
S_0 - c = Ke^{-rT} + D - p
\]
This can be read as: a short position in a call with a long position in the stock is
equivalent to a short put position plus a certain amount of cash $Ke^{-rT} + D$. 
This is a \textbf{covered call}. 

\end{document}
